% -------------------------------------------------------------------------
% WBS ID: RES-STR-RSP
% Deliverable: Structural Response and Stability
% Status: Not started
% Evidence status: Not assessed
% Review status: Not reviewed
% Last updated:
% Dependencies:
% Downstream interfaces:
% -------------------------------------------------------------------------

\section{RES-STR-RSP --- Structural Response and Stability}
\label{sec:res-str-rsp}

\subsection{Purpose and Scope}
\label{sec:res-str-rsp-purpose}

% TODO: Define what this Deliverable covers and explicitly excludes.

\subsection{Research Questions}
\label{sec:res-str-rsp-research-questions}

% TODO: State the questions that must be answered before the Deliverable
% can be accepted.

\subsection{Terminology and Definitions}
\label{sec:res-str-rsp-definitions}

% TODO: Define the technical terminology, symbols, quantities and
% classifications used in this Deliverable.

\subsection{Literature Evidence}
\label{sec:res-str-rsp-literature-evidence}

% TODO: Record evidence from peer-reviewed papers, authoritative datasets,
% standards, technical reports and primary documentation.
%
% Do not create a detached literature-summary list. Explain how each source
% supports, contradicts or limits a project decision.

\subsection{Governing Relations and Quantitative Definitions}
\label{sec:res-str-rsp-governing-relations}

% TODO: Record relevant equations, dimensionless groups, empirical models,
% extraction rules or quantitative definitions.
%
% Use align environments for displayed equations.
% Every displayed equation must have a unique \label{eq:...}.
% Do not place punctuation after displayed equations.

\subsection{Three-Dimensional Structural Equation of Motion}
\label{subsec:res_str_rsp_three_dimensional_eom}

For the active hydrodynamic baseline, the barrier is represented in
the fluid model as a fixed rigid solid. The deformable three-dimensional
solid description below is retained as a future structural-response
extension after hydrodynamic load histories have been validated.

Following spatial discretisation, the structural equation of motion is

\begin{align}
\bm{M}_{s}\ddot{\bm{q}}
+
\bm{C}_{s}\dot{\bm{q}}
+
\bm{f}_{\mathrm{int}}\!\left(\bm{q},\bm{z}\right)
&=
\bm{f}_{\mathrm{hyd}}(t)
+
\bm{f}_{\mathrm{body}}(t)
\label{eq:res_str_rsp_structural_eom}
\end{align}

where \(\bm{q}\) contains the structural displacement degrees of freedom, \(\bm{z}\) contains material-history variables, \(\bm{M}_{s}\) is the mass matrix, \(\bm{C}_{s}\) is the damping matrix and \(\bm{f}_{\mathrm{int}}\) is the nonlinear internal-force vector.

The fixed-base boundary condition is

\begin{align}
\bm{d}
&=
\bm{0}
\qquad
\text{on }\Gamma_{\mathrm{base}}
\label{eq:res_str_rsp_fixed_base_boundary}
\end{align}

A nonlinear and irregular tsunami load does not itself require nonlinear material behaviour. Structural nonlinearity arises from finite deformation, nonlinear constitutive response, contact, damage or failure. These mechanisms are deferred until a separate structural activation gate is passed.

\subsection{Updated-Lagrangian Reference Description}
\label{subsec:res_str_rsp_updated_lagrangian}

The Lagrangian formulation governs how the deforming solid is referenced during one tsunami simulation. In an Updated-Lagrangian formulation, the latest converged configuration becomes the reference for the following increment

\begin{align}
\bm{x}_{n+1}
&=
\bm{x}_{n}
+
\Delta\bm{d}_{n+1}
\label{eq:res_str_rsp_updated_lagrangian_position}
\end{align}

This differs from the barrier-class comparison loop. Each future
structural case would begin with its own undeformed design geometry,
while the Updated-Lagrangian formulation governs deformation within
that structural simulation.

If structural response is activated, candidate geometries may undergo
substantial rotation, yielding and permanent shape change. The
Updated-Lagrangian description is therefore retained as the preferred
finite-deformation reference formulation. It is compatible with
three-dimensional transient solid dynamics and incremental
constitutive updates \autocite{nemer2021}

\textbf{Deferred selected approach.} Use an Updated-Lagrangian
three-dimensional finite-deformation solid model only after the
fixed-rigid hydrodynamic load extraction and validation gate is passed.

\subsection{Material or Structural System}
\label{sec:res-str-rsp-material-structural-system}

% TODO: Complete this RES-STR Domain-specific subsection.

\subsection{Load Cases}
\label{sec:res-str-rsp-load-cases}

% TODO: Complete this RES-STR Domain-specific subsection.

\subsection{Constitutive Behaviour}
\label{sec:res-str-rsp-constitutive-behaviour}

% TODO: Complete this RES-STR Domain-specific subsection.

\subsection{Response and Stability Measures}
\label{sec:res-str-rsp-response-stability-measures}

% TODO: Complete this RES-STR Domain-specific subsection.

\subsection{Damage and Failure Modes}
\label{sec:res-str-rsp-damage-failure-modes}

% TODO: Complete this RES-STR Domain-specific subsection.

\subsection{Fluid--Structure Coupling Requirements}
\label{sec:res-str-rsp-fluid-structure-coupling-requirements}

% TODO: Complete this RES-STR Domain-specific subsection.

\subsection{Candidate Approaches}
\label{sec:res-str-rsp-candidate-approaches}

% TODO: Define the credible candidate approaches considered.

\subsection{Critical Comparison}
\label{sec:res-str-rsp-critical-comparison}

% TODO: Compare candidates using physically and project-relevant criteria.
% Do not select an approach solely on popularity or implementation ease.

\subsection{Assumptions and Applicability}
\label{sec:res-str-rsp-assumptions}

% TODO: State assumptions and the conditions under which they are valid.

\subsection{Limitations and Failure Conditions}
\label{sec:res-str-rsp-limitations}

% TODO: State where the evidence, model or selected approach ceases to be
% reliable or applicable.

\subsection{Project-Specific Conclusions}
\label{sec:res-str-rsp-conclusions}

% TODO: Convert the evidence into conclusions specific to the Studentship
% project. Do not merely restate literature findings.

\subsection{Selected Approach and Rationale}
\label{sec:res-str-rsp-selected-approach}

% TODO: Record the selected approach only after sufficient evidence exists.
% State the alternatives rejected and the reasons for rejection.

\subsection{Unresolved Decisions}
\label{sec:res-str-rsp-unresolved-decisions}

% TODO: Record unresolved questions, required evidence and decision owners.

\subsection{Requirements and Handoffs}
\label{sec:res-str-rsp-requirements}

% TODO: State requirements passed to other Research Domains, Software,
% verification activities or later execution work.

\subsection{Deliverable Completion Record}
\label{sec:res-str-rsp-completion-record}

% TODO: Record whether the Deliverable is:
%
% - Not started
% - In progress
% - Draft complete
% - Reviewed
% - Accepted
%
% Acceptance requires:
%
% - the research questions to be answered;
% - claims to be supported by appropriate evidence;
% - assumptions and limitations to be explicit;
% - the project-specific conclusion to be stated;
% - downstream requirements to be traceable;
% - unresolved decisions to be closed or formally deferred.
